% ==========================================================
% อภิธานศัพท์วิชาการ (Glossary of Agricultural Microbiology & AIoT)
% ตำรา: จุลชีววิทยาการเกษตรและเทคโนโลยี AIoT
% ระดับมาตรฐาน: รองศาสตราจารย์ (รศ.) คุณภาพ "ดีมาก"
% ==========================================================

\chapter{อภิธานศัพท์วิชาการ (Glossary)}
\label{ch:glossary}

อภิธานศัพท์นี้รวบรวมและนิยามคำศัพท์เทคนิคเฉพาะทางด้านจุลชีววิทยาการเกษตร เทคโนโลยีชีวภาพดิน ชีวเคมี พันธุศาสตร์ระดับโมเลกุล และสถาปัตยกรรมปัญญาประดิษฐ์ในสรรพสิ่ง (AIoT) โดยจัดเรียงตามลำดับอักษรภาษาอังกฤษ (A--Z) ควบคู่กับคำแปลและคำอธิบายเชิงวิชาการภาษาไทย เพื่อเป็นแหล่งอ้างอิงมาตรฐานสำหรับนักวิชาการ นักวิจัย และนักศึกษาระดับบัณฑิตศึกษา

\begin{description}[leftmargin=1.8cm, style=nextline, font=\bfseries\color{agriGreen}]

\item[1-Aminocyclopropane-1-Carboxylate (ACC) Deaminase]\index{ACC Deaminase}
เอนไซม์ในไซโทพลาซึมของแบคทีเรียส่งเสริมการเจริญเติบโตพืช (PGPR) ทำหน้าที่เร่งปฏิกิริยาสลายสารตั้งต้น ACC ให้กลายเป็นแอมโมเนียและ $\alpha$-ketobutyrate ส่งผลให้ระดับฮอร์โมนเอทิลีน ($C_2H_4$) ภายในเนื้อเยื่อพืชลดต่ำลง ช่วยบรรเทาความเครียดจากภาวะแห้งแล้ง ดินเค็ม หรือโลหะหนัก

\item[Amplicon Sequence Variant (ASV)]\index{Amplicon Sequence Variants (ASVs)}
หน่วยจำแนกสิ่งมีชีวิตในระดับความละเอียดระดับนิวคลีโอไทด์เดี่ยวที่ได้จากขั้นตอนการอนุมานลำดับเบสทางชีวสารสนเทศศาสตร์ (อาทิ อัลกอริทึม DADA2) ซึ่งแก้ไขข้อผิดพลาดในการหาลำดับเบสระดับโมเลกุลโดยไม่ต้องจัดกลุ่มตามเกณฑ์ความคล้ายคลึงแบบรวมกลุ่ม (OTU Clustering 97\%) ทำให้สามารถแยกแยะสายพันธุ์ที่มีความผันแปรเพียงหนึ่งคู่เบสได้

\item[Antagonistic Microorganisms (จุลินทรีย์ปฏิปักษ์)]\index{จุลินทรีย์ปฏิปักษ์ (Antagonistic Microorganisms)}
กลุ่มจุลินทรีย์ที่มีกลไกในการยับยั้ง ควบคุม หรือทำลายจุลินทรีย์ก่อโรคพืช ผ่านกลไกการสร้างสารปฏิชีวนะ (Antibiosis), การแย่งชิงสารอาหารและพื้นที่ (Competition), การเป็นปรสิตโดยตรง (Hyperparasitism/Mycoparasitism) หรือการกระตุ้นภูมิต้านทานของพืช (ISR/SAR) เช่น เชื้อราสกุล \textit{Trichoderma} และแบคทีเรีย \textit{Bacillus velezensis}

\item[Arbuscular Mycorrhizal Fungi (AMF)]\index{Arbuscular Mycorrhizal Fungi (AMF)}
ราเอนโดไมคอร์ไรซาในไฟลัม Glomeromycota ที่สร้างเส้นใยแทรกผ่านผนังเซลล์ของคอร์เทกซ์รากพืช แล้วขยายตัวเป็นโครงสร้างคล้ายต้นไม้ขนาดเล็ก (Arbuscules) สำหรับการแลกเปลี่ยนสารอาหาร โดยเฉพาะฟอสฟอรัสอนินทรีย์ ($H_2PO_4^-$) และน้ำ แลกกับคาร์บอนอินทรีย์ที่พืชสังเคราะห์ด้วยแสง พร้อมผลิตไกลโคโปรตีนโกลมาลิน (Glomalin) ช่วยสร้างเม็ดดินคงทน

\item[Artificial Intelligence of Things (AIoT)]\index{AIoT (Artificial Intelligence of Things)}
การหลอมรวมเทคโนโลยีปัญญาประดิษฐ์ (AI) เข้ากับโครงข่ายอินเทอร์เน็ตของสรรพสิ่ง (IoT) ทำให้โหนดฮาร์ดแวร์เซนเซอร์ภาคสนามมีความสามารถในการประมวลผล อนุมานข้อมูลเชิงลึก (Edge Inference) และตัดสินใจสั่งการควบคุมสภาวะแวดล้อมทางการเกษตรได้แบบอัตโนมัติเรียลไทม์ โดยไม่ต้องพึ่งพาเซิร์ฟเวอร์ส่วนกลางตลอดเวลา

\item[Biofloc Technology (BFT)]\index{เทคโนโลยีไบโอฟล็อก (Biofloc Technology: BFT)}
เทคโนโลยีการบำบัดน้ำและเพาะเลี้ยงสัตว์น้ำแบบหนาแน่น โดยการปรับอัตราส่วนคาร์บอนต่อไนโตรเจน (C:N Ratio $\ge 15:1$) เพื่อกระตุ้นการเพิ่มจำนวนของแบคทีเรียเฮเทอโรโทรฟิก (Heterotrophic bacteria) ซึ่งรวมกลุ่มกับสาหร่าย โปรโตซัว และเศษสารอินทรีย์ กลายเป็นกลุ่มก้อนฟล็อกชีวภาพ (Biofloc) ทำหน้าที่ดูดซับแอมโมเนียและเปลี่ยนเป็นโปรตีนเซลล์เดียวสำหรับสัตว์น้ำ

\item[Biological Nitrogen Fixation (BNF)]\index{การตรึงไนโตรเจนทางชีวภาพ (BNF)}
กระบวนการรีดิวซ์ก๊าซไนโตรเจนในบรรยากาศ ($N_2$) ให้กลายเป็นแอมโมเนีย ($NH_3$) ภายใต้สภาวะไร้ออกซิเจนโดยเอนไซม์ไนโตรจีเนส (Nitrogenase) ของจุลินทรีย์ทั้งกลุ่มที่อาศัยเป็นอิสระ (เช่น \textit{Azotobacter}) และกลุ่มที่อยู่ร่วมกับพืชแบบพึ่งพาอาศัยกัน (เช่น \textit{Rhizobium} กับพืชตระกูลถั่ว)

\item[Bray-Curtis Dissimilarity]\index{Bray-Curtis Dissimilarity}
สถิติบอกระยะห่างหรือความไม่คล้ายคลึงกัน (Dissimilarity metric) ระหว่างสองชุมชนชีวภาพ โดยคำนึงถึงทั้งการปรากฏและสัดส่วนความชุกชุมสัมพัทธ์ (Relative abundance) ของแต่ละแท็กซา มีค่าตั้งแต่ 0 (เหมือนกันทุกประการ) ถึง 1 (ไม่มีแท็กซาร่วมกันเลย) นิยมใช้ในการวิเคราะห์ Beta Diversity ทางเมตาเจโนมิกส์

\item[Carbon-to-Nitrogen Ratio (C:N Ratio)]\index{อัตราส่วน C:N (Carbon-to-Nitrogen Ratio)}
สัดส่วนมวลระหว่างคาร์บอนอินทรีย์ต่อไนโตรเจนทั้งหมดในสารอินทรีย์หรืออาหารเลี้ยงเชื้อ เป็นตัวแปรควบคุมหลักในการสลายตัวของสารอินทรีย์ การเปลี่ยนผ่านประชากรจุลินทรีย์ในการทำปุ๋ยหมัก (ค่าที่เหมาะสมเริ่มต้น $25:1$--$30:1$) และการทำงานของระบบชีวภาพ

\item[Chitinase]\index{เอนไซม์ไคทิเนส (Chitinase)}
เอนไซม์กลุ่มไฮโดรเลส ทำหน้าที่เร่งการย่อยสลายพันธะ $\beta$-1,4-glycosidic ภายในโครงสร้างสายพอลิเมอร์ของไคทิน (Chitin) ซึ่งเป็นองค์ประกอบหลักของผนังเซลล์เชื้อราก่อโรคพืชและเปลือกแมลงศัตรูพืช

\item[Clostridium Group / Methanogenesis Pathway]\index{Clostridium@\textit{Clostridium}}
กลุ่มแบคทีเรียและอาร์เคียในกระบวนการหมักแบบไร้ออกซิเจน ซึ่งเปลี่ยนสารตัวกลางกรดไขมันระเหยง่าย (VFAs) ให้กลายเป็นก๊าซมีเทน ($CH_4$) และคาร์บอนไดออกไซด์ ($CO_2$) โดยผ่านกระบวนการ Hydrolysis, Acidogenesis, Acetogenesis และ Methanogenesis

\item[Complementary DNA (cDNA)]\index{cDNA (Complementary DNA)}
สายดีเอ็นเอที่ถูกสังเคราะห์ขึ้นจากแม่แบบเอ็มอาร์เอ็นเอ (mRNA) ด้วยเอนไซม์ Reverse Transcriptase เพื่อใช้ในการตรวจวัดการแสดงออกของยีนที่กำลังทำงานในเซลล์จุลินทรีย์หรือพืช ณ เวลาใดเวลาหนึ่ง (Metatranscriptomics)

\item[Cyclic Voltammetry (CV)]\index{Cyclic Voltammetry (CV)}
เทคนิคทางเคมีไฟฟ้าที่ป้อนศักย์ไฟฟ้ากระตุ้นแบบวนลูปรูปสามเหลี่ยมไปยังขั้วไฟฟ้าทำงาน และตรวจวัดกระแสไฟฟ้าที่ตอบสนอง เพื่อวิเคราะห์ปฏิกิริยารีดอกซ์ จลนศาสตร์การถ่ายโอนอิเล็กตรอนของจุลินทรีย์สังเคราะห์กระแสไฟฟ้า หรือตรวจวัดสารบ่งชี้ทางชีวภาพ

\item[DADA2 (Divisive Amplicon Denoising Algorithm 2)]\index{DADA2 pipeline}
ไปป์ไลน์ชีวสารสนเทศศาสตร์แบบโอเพนซอร์สสำหรับการจัดเรียงลำดับนิวคลีโอไทด์จากการตรวจ 16S/18S/ITS โดยใช้แบบจำลองสถิติข้อผิดพลาดในการหาลำดับเบส (Error model) เพื่อแยกแยะความแปรผันทางพันธุกรรมที่แท้จริงออกจากสัญญาณรบกวน ส่งมอบผลลัพธ์เป็น ASV Table

\item[Dilution Rate ($D$)]\index{อัตราการเจือจาง (Dilution Rate)}
อัตราการเจือจางในระบบการเพาะเลี้ยงจุลินทรีย์แบบต่อเนื่อง (Continuous Culture หรือ Chemostat) คำนวณจากอัตราการไหลของอาหารเหลวสดเข้าสู่ถังหมัก ($F$) หารด้วยปริมาตรของเหลวในถัง ($V$) มีหน่วยเป็น $\text{h}^{-1}$ ณ สภาวะคงตัว อัตราการเจือจางจะมีค่าเท่ากับอัตราการเจริญจำเพาะ ($\mu$)

\item[Disease Severity Index (DSI)]\index{ดัชนีความรุนแรงของโรค (DSI)}
ดัชนีเชิงปริมาณที่ใช้ประเมินระดับความรุนแรงของโรคพืชในแปลงทดลอง โดยคำนวณจากผลรวมถ่วงน้ำหนักของระดับคะแนนอาการโรคพืชต่อจำนวนตัวอย่างทั้งหมด เพื่อแปลงข้อมูลคุณภาพเป็นข้อมูลปริมาณร้อยละ

\item[Edge Computing / Edge AI]\index{การตรวจจำแนกโรคพืชด้วย Edge AI}
สถาปัตยกรรมการประมวลผลข้อมูลที่ดำเนินการ ณ จุดกำเนิดข้อมูล (อุปกรณ์ปลายทาง เช่น ไมโครคอนโทรลเลอร์ ESP32, STM32 หรือบอร์ดประมวลผล Edge NPU) เพื่อลดเวลาแฝง (Latency) ประหยัดแบนด์วิดท์ และทำให้ระบบควบคุมฟาร์มสามารถอนุมานตรวจจับโรคพืชได้แม้ไม่มีสัญญาณอินเทอร์เน็ต

\item[Effector-Triggered Immunity (ETI)]\index{Effector-Triggered Immunity (ETI)}
ระบบภูมิคุ้มกันพืชระดับที่สองที่ถูกกระตุ้นเมื่อโปรตีนตัวรับภายในเซลล์พืช (NLR receptors) เข้าจับและรับรู้โปรตีนกดภูมิคุ้มกัน (Effectors) ที่เชื้อโรคปล่อยเข้ามา นำไปสู่การส่งสัญญาณตอบสนองที่รวดเร็วและรุนแรง มักเกิดปฏิกิริยาการตายของเซลล์เฉพาะที่ (Hypersensitive Response: HR)

\item[Endophyte (จุลินทรีย์เอนโดไฟต์)]\index{จุลินทรีย์เอนโดไฟต์ (Endophytes)}
แบคทีเรียหรือเชื้อราที่อาศัยอยู่ภายในเนื้อเยื่อของพืช โดยไม่ก่อให้เกิดอาการผิดปกติหรือโรคแก่พืชเจ้าบ้านในสภาวะปกติ พร้อมทั้งช่วยส่งเสริมการเจริญเติบโตหรือหลั่งสารต้านความเครียดและเชื้อโรค

\item[Glomalin-Related Soil Protein (GRSP)]\index{กลอมาลิน (Glomalin)}
สารไกลโคโปรตีนทนทานต่อการย่อยสลายทางชีวภาพ ผลิตโดยเส้นใยและสปอร์ของราไมคอร์ไรซา (AMF) ทำหน้าที่เสมือนกาวชีวภาพยึดเกาะอนุภาคดินเหนียว ดินทราย และทรายแป้งเข้าด้วยกันเป็นเม็ดดินเสถียร (Water-stable aggregates) และกักเก็บคาร์บอนอินทรีย์ในดินได้นานนับสิบปี

\item[Haldane Inhibition Kinetics]\index{แบบจำลองฮอลเดน (Haldane Kinetics)}
แบบจำลองคณิตศาสตร์จลนศาสตร์การเจริญเติบโตของจุลินทรีย์ที่ขยายขอบเขตจาก Monod เพื่ออธิบายปรากฏการณ์ที่สารตั้งต้นที่ความเข้มข้นสูงเกินไปออกฤทธิ์ยับยั้งการเจริญของเซลล์จุลินทรีย์เอง (Substrate Inhibition)

\item[Indole-3-Acetic Acid (IAA)]\index{Indole-3-acetic acid (IAA)}
ฮอร์โมนพืชในกลุ่มออกซินที่สังเคราะห์ได้โดยแบคทีเรียดินหลายชนิดผ่านวิถีชีวเคมีที่พึ่งพากรดอะมิโน L-Tryptophan ช่วยกระตุ้นการยืดตัวของเซลล์ราก การแตกแขนงของรากฝอย และการเพิ่มพื้นที่ผิวในการดูดซับธาตุอาหาร

\item[Induced Systemic Resistance (ISR)]\index{Induced Systemic Resistance (ISR)}
กลไกการต้านทานโรคทั้งต้นของพืชที่ถูกกระตุ้นโดยแบคทีเรียหรือราที่เป็นมิตรในดิน (เช่น PGPR หรือ \textit{Trichoderma}) โดยไม่ทำให้เกิดอาการตายเฉพาะที่ สัญญาณภูมิคุ้มกันนี้ส่งผ่านวิถีของกรดจัสโมนิก (Jasmonic acid) และเอทิลีน (Ethylene)

\item[Keystone Taxa (สายพันธุ์หลักในระบบนิเวศ)]\index{Keystone Hub Taxa}
กลุ่มจุลินทรีย์ที่มีบทบาทสำคัญอย่างยิ่งยวดในการค้ำจุนโครงสร้างและเสถียรภาพการทำงานของชุมชนจุลินทรีย์ทั้งหมด หากแท็กซาเหล่านี้ถูกกำจัดหรือสูญเสียไป เครือข่ายการแลกเปลี่ยนสารอาหารและการทำงานของไมโครไบโอมจะล่มสลายลง แม้จะมีสัดส่วนความชุกชุมสัมพัทธ์ในดินไม่สูงก็ตาม

\item[LoRaWAN (Long Range Wide Area Network)]\index{LoRaWAN (Long Range Wide Area Network)}
โพรโทคอลเครือข่ายสื่อสารไร้สายพลังงานต่ำระยะไกล ทำงานบนคลื่นความถี่คลื่นวิทยุย่าน Sub-GHz (920--925 MHz สำหรับประเทศไทย) เหมาะสำหรับการส่งข้อมูลโทรมาตรจากเซนเซอร์ภาคสนามที่อยู่ห่างออกไปหลายกิโลเมตรกลับสู่เกตเวย์

\item[Metagenomics (เมตาเจโนมิกส์)]\index{เมตาเจโนมิกส์ (Metagenomics)}
การศึกษาสารพันธุกรรมทั้งหมดที่สกัดได้โดยตรงจากตัวอย่างสิ่งแวดล้อม (อาทิ ดิน น้ำ หรือทางเดินอาหารสัตว์) โดยไม่ต้องผ่านการเพาะเลี้ยงเชื้อในห้องปฏิบัติการ ทำให้สามารถตรวจพบและจำแนกจุลินทรีย์ที่ไม่สามารถเพาะเลี้ยงได้ ซึ่งคิดเป็นกว่า 99\% ของจุลินทรีย์ในธรรมชาติ

\item[Methanosaeta / Methanosarcina]\index{Methanosaeta@\textit{Methanosaeta}}
จีนัสของอาร์เคียกลุ่มอะซีโทโทรฟิก (Acetotrophic methanogens) ที่มีความสามารถเฉพาะทางในการเปลี่ยนกรดแอซีติก ($CH_3COOH$) ให้กลายเป็นก๊าซมีเทนและคาร์บอนไดออกไซด์ เป็นองค์ประกอบชี้วัดความเสถียรของระบบถังหมักก๊าซชีวภาพ

\item[Microbial Carbon Pump (MCP)]\index{Microbial Carbon Pump (MCP)}
แนวคิดทางนิเวศวิทยาชีวเคมีที่อธิบายกระบวนการที่จุลินทรีย์ดินเปลี่ยนรูปสารประกอบคาร์บอนอินทรีย์ที่ย่อยสลายง่าย (Labile Organic Carbon) ให้กลายเป็นมวลชีวภาพของจุลินทรีย์ และเมื่อเซลล์ตายจะทิ้งซากเซลล์ (Necromass) ซึ่งมีโครงสร้างต้านทานการย่อยสลาย กลายเป็นคาร์บอนอินทรีย์เสถียรในดิน (Recalcitrant Soil Organic Matter)

\item[Microencapsulation (การห่อหุ้มระดับไมโคร)]\index{การห่อหุ้มเซลล์ระดับไมโคร (Microencapsulation)}
กระบวนการกักเก็บเซลล์จุลินทรีย์มีชีวิตหรือสารชีวภาพไว้ภายในเม็ดแคปซูลขนาดเส้นผ่านศูนย์กลางระดับไมโครเมตร โดยใช้วัสดุพอลิเมอร์ชีวภาพ (เช่น แคลเซียมอัลจิเนตเคลือบไคโตซาน) เพื่อปกป้องจุลินทรีย์จากอุณหภูมิ แสงแดด สภาวะแห้งแล้ง และยืดอายุการเก็บรักษาชีวภัณฑ์

\item[Modbus RTU Protocol]\index{โพรโทคอล Modbus RTU}
โพรโทคอลสื่อสารแบบ Master-Slave ระดับอุตสาหกรรมที่ทำงานบนมาตรฐานสัญญาณไฟฟ้าเชิงกายภาพ RS485 นิยมใช้เชื่อมต่อเซนเซอร์ตรวจวัดฟิสิกส์เคมีในดิน (เช่น ค่า pH, NPK, ความชื้น, EC, อุณหภูมิ) กับตัวควบคุมไมโครคอนโทรลเลอร์

\item[Modified Gompertz Model]\index{แบบจำลอง Modified Gompertz (Modified Gompertz Model)}
แบบจำลองคณิตศาสตร์เชิงประจักษ์แบบไม่เป็นเชิงเส้นที่นิยมใช้อย่างกว้างขวางเพื่อทำนายศักยภาพการผลิตก๊าซชีวภาพสะสม อัตราการผลิตสูงสุด และระยะหน่วง (Lag phase time) ในระบบการย่อยสลายแบบไร้ออกซิเจน

\item[Monod Kinetic Model]\index{แบบจำลองโมโนด์ (Monod Model)}
แบบจำลองจลนศาสตร์พื้นฐานที่ใช้อธิบายความสัมพันธ์ระหว่างอัตราการเจริญจำเพาะของจุลินทรีย์ ($\mu$) กับความเข้มข้นของสารอาหารจำกัดการเจริญ ($S$) ผ่านค่าพารามิเตอร์อัตราการเจริญสูงสุด ($\mu_{\max}$) และค่าคงที่การสังเคราะห์ ($K_s$)

\item[Mycoparasitism (ภาวะปรสิตของรา)]\index{Mycoparasitism (ภาวะปรสิตของรา)}
กลไกที่เชื้อราปฏิปักษ์ (เช่น \textit{Trichoderma}) เข้าจู่โจม สัมผัส พันรัด และแทรกโครงสร้างเจาะ (Appressorium-like structures) เข้าสู่เส้นใยของเชื้อราก่อโรคพืช พร้อมทั้งหลั่งเอนไซม์ย่อยสลายผนังเซลล์ (Cell wall-degrading enzymes: CWDEs) เช่น ไคทิเนสและกลูคาเนสเพื่อกินเนื้อเยื่อของราก่อโรค

\item[Nitrification (ไนตริฟิเคชัน)]\index{ไนตริฟิเคชัน (Nitrification)}
กระบวนการออกซิไดซ์สารประกอบแอมโมเนียหรือแอมโมเนียม ($NH_4^+$) ให้กลายเป็นไนไตรต์ ($NO_2^-$) โดยแบคทีเรียกลุ่ม AOB (เช่น \textit{Nitrosomonas}) จากนั้นออกซิไดซ์ต่อเป็นไนเตรต ($NO_3^-$) โดยแบคทีเรียกลุ่ม NOB (เช่น \textit{Nitrobacter}) ภายใต้สภาวะที่มีออกซิเจน

\item[Nitrogenase Enzyme Complex]\index{เอนไซม์ไนโตรจีเนส (Nitrogenase Complex)}
คอมเพล็กซ์เอนไซม์ชีวภาพประกอบด้วยสองหน่วยย่อย ได้แก่ โปรตีนเหล็ก (Fe protein) และโปรตีนโมลิบดีนัม-เหล็ก (MoFe protein) ทำหน้าที่เร่งปฏิกิริยาการรีดิวซ์ก๊าซไนโตรเจน โดยเอนไซม์นี้จะถูกทำลายสภาพอย่างรวดเร็วเมื่อสัมผัสกับก๊าซออกซิเจนอิสระ

\item[Operational Taxonomic Unit (OTU)]\index{OTU (Operational Taxonomic Unit)}
กลุ่มของลำดับนิวคลีโอไทด์ที่ถูกจัดกลุ่มรวมกันตามเกณฑ์ความเหมือนที่กำหนด (โดยทั่วไปคือ $97\%$ สำหรับ 16S rRNA gene) เพื่อใช้แทนหน่วยจำแนกชนิดจุลินทรีย์เชิงปฏิบัติการในยุคแรกของชีวสารสนเทศศาสตร์

\item[Pattern-Triggered Immunity (PTI)]\index{PAMP-Triggered Immunity (PTI)}
กลไกภูมิคุ้มกันพืชขั้นพื้นฐานที่เกิดขึ้นเมื่อตัวรับบนผิวเยื่อหุ้มเซลล์พืช (Pattern Recognition Receptors: PRRs) ตรวจจับโมเลกุลจำเพาะของเชื้อก่อโรค (PAMPs เช่น แฟลกเจลลิน หรือ ไคทิน) นำไปสู่การกระตุ้นการสร้างสารอนุมูลอิสระ (ROS) และการสะสมแคลโลสที่ผนังเซลล์

\item[Phosphate-Solubilizing Bacteria (PSB)]\index{แบคทีเรียละลายฟอสเฟต (PSB)}
กลุ่มแบคทีเรียดินที่มีความสามารถในการเปลี่ยนฟอสฟอรัสอนินทรีย์ที่ไม่ละลายน้ำ (เช่น แคลเซียมฟอสเฟต เหล็กฟอสเฟต หรืออะลูมิเนียมฟอสเฟต) ให้อยู่ในรูปที่พืชดูดซึมได้ ($H_2PO_4^-$ หรือ $HPO_4^{2-}$) ผ่านการหลั่งกรดอินทรีย์โมเลกุลต่ำ (เช่น กรดซิตริก กรดกลูโคนิก) หรือการหลั่งเอนไซม์ฟอสฟาเทส

\item[Plant Growth-Promoting Rhizobacteria (PGPR)]\index{PGPR (Plant Growth-Promoting Rhizobacteria)}
กลุ่มแบคทีเรียที่มีประโยชน์ซึ่งอาศัยอยู่อย่างหนาแน่นในบริเวณไรโซสเฟียร์และส่งเสริมการเจริญเติบโตของพืชผ่านกลไกทางตรง (เช่น การตรึงไนโตรเจน การละลายฟอสฟอรัส การสังเคราะห์ไฟโตฮอร์โมน) และกลไกทางอ้อม (เช่น การยับยั้งเชื้อก่อโรคและการกระตุ้นภูมิต้านทานของพืช)

\item[Polyhydroxyalkanoates (PHAs)]\index{พอลิไฮดรอกซีอัลคาโนเอต (PHAs)}
พอลิเมอร์ชีวภาพกลุ่มพอลิเอสเทอร์ที่จุลินทรีย์หลายชนิดสังเคราะห์และสะสมไว้ภายในเซลล์เป็นแหล่งพลังงานสำรองเมื่อเกิดสภาวะสารอาหารอื่นจำกัด สามารถนำมาสกัดใช้เป็นพลาสติกชีวภาพที่ย่อยสลายได้ตามธรรมชาติ 100\%

\item[Quorum Sensing (QS)]\index{การรับรู้จำนวนประชากร (Quorum Sensing)}
ระบบการสื่อสารระหว่างเซลล์จุลินทรีย์ที่ขึ้นกับความหนาแน่นของประชากร โดยจุลินทรีย์จะหลั่งสารส่งสัญญาณโมเลกุลขนาดเล็ก (Autoinducers) เช่น $N$-acyl homoserine lactones (AHLs) เมื่อความเข้มข้นถึงจุดวิกฤติ จะกระตุ้นการแสดงออกของยีนที่ควบคุมพฤติกรรมกลุ่ม เช่น การสร้างไบโอฟิล์ม การเรืองแสง หรือการปลดปล่อยสารพิษ

\item[Random Forest Regressor (RFR)]\index{Random Forest Regression}
อัลกอริทึมการเรียนรู้ของเครื่องแบบ Ensemble Learning ที่สร้างต้นไม้ตัดสินใจ (Decision Trees) จำนวนมากเพื่อทำนายค่าตัวแปรตามต่อเนื่อง โดยการสุ่มเลือกชุดข้อมูลฝึกและตัวแปรย่อย เหมาะสำหรับการทำนายความหนาแน่นจุลินทรีย์หรือสภาวะเคมีในแปลงเกษตรที่มีตัวแปรร่วมจำนวนมาก

\item[Rhizosphere Competence]\index{Rhizosphere Competence}
ความสามารถของจุลินทรีย์ที่นำเข้าสู่ดินในการมีชีวิตรอด เติบโต แย่งชิงพื้นที่ผิวรากพืช และแพร่กระจายตัวไปตามการยืดขยายของระบบราก ภายใต้แรงกดดันและการแข่งขันจากจุลินทรีย์ประจำถิ่นเดิม

\item[Shannon-Wiener Diversity Index ($H'$)]\index{Shannon Diversity Index ($H'$)}
ดัชนีทางสถิติที่ใช้วัดความหลากหลายของสปีชีส์ (Alpha Diversity) ในชุมชนสิ่งมีชีวิต โดยคำนวณจากทั้งจำนวนชนิดที่มีอยู่ (Species Richness) และความสม่ำเสมอในการกระจายตัวของแต่ละชนิด (Evenness)

\item[Siderophores]\index{ไซเดอโรฟอร์ (Siderophores)}
สารประกอบอินทรีย์น้ำหนักโมเลกุลต่ำที่มีความจำเพาะสูงมากในการจับกับเหล็กประจุสาม ($Fe^{3+}$) ที่แบคทีเรียและราสังเคราะห์และหลั่งออกสู่ภายนอกเซลล์ เพื่อดึงดูดและขนส่งธาตุเหล็กกลับเข้าสู่เซลล์ ทั้งยังช่วยแย่งชิงธาตุเหล็กจากเชื้อโรคพืชในบริเวณใกล้เคียง

\item[Solid-State Fermentation (SSF)]\index{การหมักแบบอาหารแข็ง (SSF)}
กระบวนการหมักทางชีวภาพที่จุลินทรีย์เจริญเติบโตบนอนุภาคของแข็งที่มีความชื้นจำกัดและไม่มีน้ำอิสระไหลเวียน เหมาะสำหรับการผลิตสปอร์ของเชื้อราปฏิปักษ์ เช่น \textit{Trichoderma} และ \textit{Beauveria} บนเมล็ดธัญพืช

\item[Submerged Fermentation (SmF)]\index{การหมักในอาหารเหลว (SmF)}
กระบวนการหมักในสารละลายอาหารเหลวที่มีการกวนผสมและเติมอากาศอย่างต่อเนื่องในถังปฏิกรณ์ชีวภาพ (Bioreactor) เหมาะสำหรับการผลิตมวลชีวภาพแบคทีเรีย เอนไซม์ หรือเมแทบอไลต์ทุติยภูมิในระดับอุตสาหกรรม

\item[Synthetic Microbial Communities (SynComs)]\index{กลุ่มจุลินทรีย์สังเคราะห์ (SynComs)}
กลุ่มจุลินทรีย์สังเคราะห์ที่ได้รับการออกแบบ ผสมผสาน และประกอบขึ้นจากสายพันธุ์จุลินทรีย์บริสุทธิ์ที่มีบทบาทหน้าที่เสริมฤทธิ์กัน (Synergistic function) ตามผลการวิเคราะห์เครือข่ายเมตาเจโนมิกส์ เพื่อสร้างชุมชนจุลินทรีย์ที่มีเสถียรภาพและประสิทธิภาพสูงกว่าสายพันธุ์เดี่ยว

\item[Systemic Acquired Resistance (SAR)]\index{Systemic Acquired Resistance (SAR)}
กลไกการต้านทานโรคทั่วทั้งต้นของพืชที่ถูกกระตุ้นขึ้นหลังจากการติดเชื้อโรคพืชเฉพาะที่ โดยมีการส่งสัญญาณเตือนผ่านวิถีของกรดซาลิไซลิก (Salicylic acid) นำไปสู่การสังเคราะห์โปรตีนต้านโรคพืช (PR proteins) ทั่วทั้งระบบ

\item[Volatile Organic Compounds (VOCs)]\index{สารประกอบอินทรีย์ระเหยง่าย (VOCs)}
สารประกอบอินทรีย์ระเหยง่ายที่มีน้ำหนักโมเลกุลต่ำและความดันไอสูง ผลิตโดยจุลินทรีย์ดิน สามารถแพร่กระจายผ่านรูพรุนในดินและบรรยากาศ ทำหน้าที่เป็นสัญญาณกระตุ้นการเจริญของพืช ยับยั้งเชื้อรา หรือดึงดูดสิ่งมีชีวิตที่เป็นประโยชน์

\item[Volumetric Oxygen Transfer Coefficient ($k_L a$)]\index{สัมประสิทธิ์การถ่ายโอนออกซิเจน ($k_L a$)}
สัมประสิทธิ์การถ่ายโอนออกซิเจนเชิงปริมาตรในถังหมักชีวภาพ มีหน่วยเป็น $\text{h}^{-1}$ หรือ $\text{s}^{-1}$ เป็นตัวแปรวิศวกรรมที่บ่งบอกถึงประสิทธิภาพในการละลายของก๊าซออกซิเจนจากฟองก๊าซสู่สารอาหารเลี้ยงเชื้อเหลว

\item[Winogradsky Column]\index{Winogradsky Column}
อุปกรณ์จำลองระบบนิเวศจุลินทรีย์แบบปิดในหลอดแก้ว พัฒนาโดย Sergei Winogradsky เพื่อศึกษาความหลากหลายและการจัดชั้นทางนิเวศวิทยาของจุลินทรีย์ตามเกรเดียนต์ของความเข้มข้นออกซิเจนและซัลไฟด์

\item[Yield Coefficient ($Y_{X/S}$)]\index{สัมประสิทธิ์ผลได้ชีวมวล ($Y_{X/S}$)}
ค่าสัมประสิทธิ์ผลได้ของการเจริญ บ่งบอกถึงมวลชีวภาพของจุลินทรีย์ที่สร้างขึ้น ($\Delta X$) ต่อปริมาณสารตั้งต้นคาร์บอนที่ถูกใช้ไป ($\Delta S$) มีหน่วยเป็น $\text{g biomass/g substrate}$

\item[YOLOv8-Nano (You Only Look Once)]\index{YOLOv8-Nano Deep Learning}
สถาปัตยกรรมโครงข่ายประสาทเทียมแบบคอนโวลูชันสำหรับการตรวจจับวัตถุแบบเรียลไทม์ (Object Detection) ขนาดกะทัดรัดพิเศษ ออกแบบมาให้ทำงานบนอุปกรณ์ประมวลผลขนาดเล็ก (Edge Embedded AI) เพื่อตรวจจับอาการโรคพืช รอยโรคบนใบ และศัตรูพืชได้อย่างแม่นยำและรวดเร็ว

\item[Zigzag Model of Plant Immunity]\index{Zigzag Model of Plant Immunity}
แบบจำลองเชิงแนวคิดที่อธิบายวิวัฒนาการร่วมระหว่างระบบภูมิคุ้มกันพืชกับเชื้อก่อโรคพืช แบ่งออกเป็น 4 เฟสหลัก ได้แก่ การรับรู้ PAMPs (PTI), การปล่อย Effectors ยับยั้งภูมิคุ้มกัน (ETS), การตรวจจับ Effector โดยพืช (ETI) และการกลายพันธุ์ของ Effector เพื่อหลบเลี่ยงการตรวจจับรอบใหม่

\end{description}
